\subsection{LOCALLY COMPACT HOMOGENEOUS SPACES}

PROPOSITION 13. Let G be a locally compact group and let H be a closed subgroup of G. Then the homogeneous space G/H is locally compact and paracompact.

Since G/H is Hausdorff (§ 2, no. 5, Proposition 13) it is locally compact, by Proposition 9 of no. 5 applied to H operating on G on the right. Thus it remains to show that G/H is paracompact. Let V be a symmetric compact neighbourhood of e in G, and let \(G_{0}=V^{\infty}\) be the subgroup of G generated by V. \(G_{0}\) is open (§ 2, no. 1, Corollary to Proposition 4) and operates continuously on G/H (§ 2, no. 5, Proposition 12). If we can show that each of the orbits \(G_{0}.z\) ( \(z\in G/H\) ) is an open subset of G/H and a countable union of compact sets, then it will follow that G/H is the topological sum of the distinct orbits \(G_{0}.z\) and is therefore paracompact (Chapter I, § 9, no. 10, Theorem 5). That \(\mathbf{G}_0.z\) is open in G/H follows from the facts that \(\mathbf{G}_0\) is open in G and the equivalence relation \(x^{-1}y\in H\) is open (§ 2, no. 4, Lemma 2). On the other hand, \(\mathbf{G}_0.z\) is the union of the sets \(\mathbf{V}^n.z\) ( \(n\geqslant 1\) ), and since \(\mathbf{V}^n\) is compact in G and G/H is Hausdorff, \(\mathbf{V}^n.z\) is compact. This completes the proof.

PROPOSITION 14. In a locally compact group G, the identity component C is the intersection of the open subgroups of G.

C is a closed normal subgroup of G (§ 2, no. 2, Proposition 7), and hence G/C is locally compact (Proposition 13) and totally disconnected (Chapter I, § 11, no. 5, Proposition 9). Since the inverse image of an open subgroup of G/C, under the canonical mapping of G onto G/C, is an open subgroup of G containing C, we see that we can restrict ourselves to proving the proposition for the group G/C. In other words, we may assume that G is totally disconnected. We know then (Chapter II, § 4, no. 4, Corollary to Proposition 6) that every compact neighbourhood V of e contains a neighbourhood U of e which is both open and closed. Since U is compact and \(B = \complement U\) is closed, there is a symmetric open neighbourhood W of e such that W ⊂ U and UW ∩ BW = ∅ (§ 3, no. 1 and Chapter II, § 4, no. 3, Proposition 4), and a fortiori UW ⊂ U. By induction on n it follows that \(W^n \subset U\) for every integer \(n > 0\). Hence the subgroup \(W^\infty = \bigcup_{n > 0} W^n\), generated by W, is contained in U; but \(W^\infty\) is open in G (§ 2, no. 1, Corollary to Proposition 4). This completes the proof.

We have also proved :

COROLLARY 1. If G is a totally disconnected locally compact group, then every neighbourhood of e in G contains an open subgroup of G.

COROLLARY 2. A locally compact group is connected if it is generated by each neighbourhood of the identity element.

COROLLARY 3. Let G be a locally compact group, let H be a closed subgroup of G, and let \(\varphi\) be the canonical mapping of G onto G/H. Then the components of G/H are the closures of the images under \(\varphi\) of the components of G.

Let C be the identity component of G. The components of G are the sets sC, where \(s \in G\) (§ 2, no. 2, Proposition 7); \(\varphi(sC)\) is clearly connected, hence so is \(\overline{\varphi(sC)}\) (Chapter I, § 11, no. 1, Proposition 1). But \(\varphi(sC) = \varphi(sCH)\) , and since sCH is saturated with respect to the equivalence relation defined by H, and since this equivalence relation is open

(§ 2, no. 4, Lemma 2), we have \(\overline{\varphi(s\mathbf{CH})} = \varphi(\overline{s\mathbf{CH}}) = \varphi(s.\overline{\mathbf{CH}})\) (Chapter I, § 5, no. 3, Proposition 7).

Put L = CH. L is a closed subgroup of G which contains C and H; hence to prove that the sets \(\varphi(s.L) = s.\varphi(L)\) are the components of G/H it is enough to show that the quotient space of G/H by the equivalence relation whose classes are the sets \(s.\varphi(L)\) is totally disconnected. Now this quotient space is homeomorphic to the homogeneous space G/L (Chapter I, § 3, no. 4, Proposition 7); we are thus reduced to proving that when \(C \subset H\) , G/H is totally disconnected. Since G/H may be identified with \((\mathrm{G/C})/(\mathrm{H/C})\) (§ 2, no. 7, Proposition 22), we may even assume that G itself is totally disconnected. Every neighbourhood of \(\varphi(e)\) in G/H contains a neighbourhood of the form \(\varphi(V)\) , where V is a neighbourhood of e in G, and therefore (Corollary 1) contains a neighbourhood of the form \(\varphi(K)\) , where K is a compact open subgroup of G. \(\varphi(K)\) is therefore both open and closed in G/H, and this shows that the component of \(\varphi(e)\) in G/H consists of \(\varphi(e)\) alone. By translation, the same is true for the component of every point of G/H, and the corollary is proved.

